\section{Herleitung aus der Fourier-Transformation}
\label{laplace:subsection:herleitung}

Um die Natur der inversen Laplace-Transformation zu verstehen, lohnt sich ein kurzer Blick auf ihren Ursprung: die Fourier-Transformation.
Die Fourier-Transformation zerlegt ein Signal in seine Frequenzanteile, setzt jedoch voraus, dass das Signal absolut integrierbar ist.
Viele physikalisch relevante Funktionen, wie etwa exponentiell wachsende Signale, erfüllen diese Bedingung nicht; ihr Fourier-Integral divergiert.

Um dieses Konvergenzproblem zu lösen, bedient man sich eines eleganten Tricks.
Man multipliziert die ursprüngliche Funktion $f(t)$ mit einem exponentiellen Dämpfungsfaktor $e^{-\gamma t}$, wobei $\gamma$ eine reelle Konstante ist.
Wählt man $\gamma$ gross genug, wird das Wachstum von $f(t)$ unterdrückt.
Wir definieren also eine gedämpfte Hilfsfunktion
\begin{equation}
    g(t) = f(t) e^{-\gamma t}.
    \label{laplace:eq:hilfsfunktion}
\end{equation}
Für diese gedämpfte Funktion existiert nun die Fourier-Transformation $G(\omega) = \mathcal{F}\{g(t)\}$:
\begin{equation}
    G(\omega) = \int_{-\infty}^{\infty} \left( f(t) e^{-\gamma t} \right) e^{-i\omega t} \, dt = \int_{-\infty}^{\infty} f(t) e^{-(\gamma + i\omega)t} \, dt.
    \label{laplace:eq:fourier_g}
\end{equation}
An dieser Stelle kommt uns die Physik zu Hilfe.
Die von uns betrachteten Vorgänge beginnen zu einem definierten Zeitpunkt, den wir als $t = 0$ wählen; für $t < 0$ setzen wir $f(t) = 0$ und nennen ein solches Signal \emph{kausal}.
Für kausale Funktionen verschwindet der Integrand in Gleichung~\eqref{laplace:eq:fourier_g} für alle negativen Zeiten, sodass sich das Integral effektiv nur noch von $0$ bis $\infty$ erstreckt.
Führen wir nun die komplexe Variable $s = \gamma + i\omega$ ein, so entspricht dieses Integral exakt der Definition der Laplace-Transformation von $f(t)$.
Es gilt also $G(\omega) = F(s)$.

Um nun $f(t)$ zurückzugewinnen, wenden wir die inverse Fourier-Transformation auf $G(\omega)$ an:
\begin{equation}
    g(t) = \frac{1}{2\pi} \int_{-\infty}^{\infty} G(\omega) e^{i\omega t} \, d\omega.
    \label{laplace:eq:inv_fourier_g}
\end{equation}
Setzen wir in diese Gleichung $g(t) = f(t) e^{-\gamma t}$ wieder ein und multiplizieren beide Seiten mit $e^{\gamma t}$, erhalten wir
\begin{equation}
    f(t) = \frac{e^{\gamma t}}{2\pi} \int_{-\infty}^{\infty} G(\omega) e^{i\omega t} \, d\omega 
         = \frac{1}{2\pi} \int_{-\infty}^{\infty} G(\omega) e^{(\gamma + i\omega)t} \, d\omega.
    \label{laplace:eq:rueck_1}
\end{equation}
Im letzten Schritt ersetzen wir die Integrationsvariable $\omega$ wieder durch unsere komplexe Variable $s$.
Mit $s = \gamma + i\omega$ und konstantem $\gamma$ folgt für das Differential $ds = i \, d\omega$, also $d\omega = \frac{ds}{i}$.
Da $\omega$ von $-\infty$ bis $\infty$ läuft, bewegt sich $s$ entlang einer vertikalen Geraden in der komplexen Ebene von $\gamma - i\infty$ bis $\gamma + i\infty$.
Setzen wir dies sowie $G(\omega) = F(s)$ in Gleichung~\eqref{laplace:eq:rueck_1} ein, so erhalten wir die fundamentale Formel für die inverse Laplace-Transformation:
\begin{equation}
    f(t) = \frac{1}{2\pi i} \int_{\gamma - i\infty}^{\gamma + i\infty} F(s) e^{st} \, ds.
    \label{laplace:eq:bromwich_integral}
\end{equation}

Dieses komplexe Linienintegral wird als \emph{Bromwich-Integral} bezeichnet.
Abbildung~\ref{laplace:fig:bromwich} veranschaulicht die geometrische Situation in der komplexen Ebene.
Die vertikale Integrationsgerade $\text{Re}(s) = \gamma$ wird dabei stets so gewählt, dass sie rechts von allen Singularitäten der Bildfunktion $F(s)$ liegt.
Man nennt den minimalen Wert, für den das Integral noch konvergiert, die Konvergenzabszisse $\gamma_0$.
Für alle $\gamma > \gamma_0$ liefert das Integral die eindeutige Lösung.

\begin{figure}
\centering
\scalebox{0.85}{
\begin{tikzpicture}

% ==============================
% 1. FARBEN DEFINIEREN
% ==============================
\definecolor{plotblue}{RGB}{0, 0, 200}
\definecolor{plotred}{RGB}{220, 30, 30}
\definecolor{plotgray}{RGB}{140, 140, 140}
\definecolor{plotbg}{RGB}{244, 244, 252}

% ==============================
% 2. HINTERGRUND (Konvergenzbereich)
% ==============================
\fill[plotbg] (1.5, -4.8) rectangle (5.0, 4.8);

% ==============================
% 3. ACHSEN
% ==============================
\draw[thick, -{Stealth[length=3.5mm, width=2.5mm]}] (-5.2, 0) -- (5.5, 0) node[right] {\textbf{Re}$(s)$};
\draw[thick, -{Stealth[length=3.5mm, width=2.5mm]}] (0, -5.2) -- (0, 5.2) node[above] {\textbf{Im}$(s)$};

% ==============================
% 4. LINIEN & BESCHRIFTUNGEN
% ==============================
% Konvergenzabszisse (\gamma_0)
\draw[thick, plotgray, dashed] (1.5, -4.8) -- (1.5, 4.8);
\draw[thick] (1.5, 0.15) -- (1.5, -0.15) node[below right, inner sep=2pt] {$\boldsymbol{\gamma_0}$};

% Bromwich-Pfad (\gamma)
\draw[thick, plotblue, ->-=0.55] (2.5, -4.8) -- (2.5, 4.8);
\draw[thick] (2.5, 0.15) -- (2.5, -0.15) node[below right, inner sep=2pt] {$\boldsymbol{\gamma}$};

% Text im Konvergenzbereich
\node[plotblue, align=center, font=\small] at (3.75, 3) {Konvergenz-\\bereich};

% ==============================
% 5. POLE (Rote Kreuze)
% ==============================
\newcommand{\pole}[2]{
    \draw[thick, plotred, line cap=round] (#1-0.12,#2-0.12) -- (#1+0.12,#2+0.12);
    \draw[thick, plotred, line cap=round] (#1-0.12,#2+0.12) -- (#1+0.12,#2-0.12);
}

% Pole platzieren
\pole{1.5}{0}
\pole{0}{0}
\pole{0}{2}
\pole{0}{-2}
\pole{-1.2}{1.5}
\pole{-1.2}{-1.5}
\pole{-2}{0}
\pole{-3}{0}

% ==============================
% 6. LEGENDE
% ==============================
\node[draw, thick, rounded corners=3pt, fill=white, inner xsep=5pt, inner ysep=4pt, anchor=north west] at (-5.0, 4.8) {
    \renewcommand{\arraystretch}{1.2}
    \begin{tabular}{@{}c l@{}}
        
        % Symbol 1: Polstelle
        \tikz[baseline=-0.6ex]{
            \draw[thick, plotred, line cap=round] (-0.1,-0.1) -- (0.1,0.1) (-0.1,0.1) -- (0.1,-0.1);
        } & Polstelle \\
        
        % Symbol 2: Konvergenzabszisse (etwas verlängert auf -0.3 bis 0.3)
        \tikz[baseline=-0.6ex]{
            \draw[thick, plotgray, dashed] (-0.3, 0) -- (0.3, 0);
        } & Konvergenzabszisse \\
        
        % Symbol 3: Bromwich-Pfad (Mit Dekoration für perfekten Strich)
        \tikz[baseline=-0.6ex]{
            \draw[thick, plotblue, ->-=0.6] (-0.3, 0) -- (0.3, 0); % Durchgehender Strich
        } & Bromwich-Pfad \\
        
    \end{tabular}
};

\end{tikzpicture}
}
\caption{Die Bromwich-Kontur in der komplexen $s$-Ebene. Die Integrationsgerade (blau) liegt rechts von der Konvergenzabszisse $\gamma_0$ und allen Singularitäten (rote Kreuze).}
\label{laplace:fig:bromwich}
\end{figure}

